\subsection{Wrapped and Laid}\label{sec:wrapped}

Read two verses of St Luke side by side. The first is from the night of the
Nativity:

\begin{quote}
And she brought forth her first born son and wrapped him up in swaddling clothes
and laid him in a manger: because there was no room for them in the inn.\par
{\smallskip\noindent\small ---Luke 2:7\par}
\end{quote}

The second is from the day of the Crucifixion, when Joseph of Arimathea has
begged Pilate for the body:

\begin{quote}
And taking him down, he wrapped him in fine linen and laid him in a sepulchre
that was hewed in stone, wherein never yet any man had been laid.\par
{\smallskip\noindent\small ---Luke 23:53\par}
\end{quote}

Wrapped and laid, at the beginning and at the end. The Latin Vulgate uses the
same verb for both wrappings: Mary \latin{pannis eum involvit},
Joseph \latin{involvit sindone}. The swaddling clothes and the shroud are one
gesture made twice, by a mother over a newborn and by a disciple over a corpse.
The first was the sign the angel gave the shepherds: ``You shall find the
infant wrapped in swaddling clothes and laid in a manger.''

\subsubsection*{Bound hand and foot}

A child in swaddling clothes is bound. St Bede noticed it: ``He by Whom all
things are made, is folded both hands and feet, that our hands might be raised
up for every good work, and our feet directed in the way of peace.'' In his own
Latin the phrase is sharper, \latin{manus pedesque cunis astringitur}: he is
bound hand and foot in the cradle.

He will be bound again. In the garden ``the band and the tribune and the
servants of the Jews took Jesus and bound him''; in the morning ``they brought
him bound and delivered him to Pontius Pilate the governor'', the same Pilate
whose name stands in the Gospel of the Fourth Sunday.
His hands and feet, bound in the cradle, will be dug through, as the Psalm says
whose first words he cried from the Cross: ``They have dug my hands and
feet.'' And after his
death the binding is done a last time, with love: ``They took therefore the
body of Jesus and bound it in linen cloths, with the spices, as the manner of
the Jews is to bury.''

He enters the world bound and is laid in a feeding trough. He leaves it bound,
condemned and executed, and is laid in a tomb cut from the rock.

\subsubsection*{Born outside, dying outside}

There was no room for them in the inn. St Gregory the Great, in the Christmas
homily read at Matins, sees a mystery in it:

\begin{quote}
He was born, not in His Mother's house, but away from home. And this is a
mystery, showing that this our mortality into which He was born was not the home
of Him Who is begotten of the Father before the worlds.\par
{\smallskip\noindent\small ---St Gregory the Great, Homily 8 on the Gospels\par}
\end{quote}

He was born outside, and he died outside: ``Jesus also, that he might sanctify
the people by his own blood, suffered without the gate.'' St Bede takes
Gregory's thought further: he is born ``in an inn and by the way side, because
through the mystery of the incarnation He was made the way by which to bring us
to our country''.

The gifts of the wise men point the same way. St Matthew says that,
``opening their treasures, they offered him gifts; gold, frankincense, and
myrrh.'' At the burial Nicodemus came ``bringing a mixture of myrrh and aloes,
about an hundred pound weight''. St Gregory Nazianzen tells his hearers at
Christmas to run with the star and bear their gifts with the Magi, ``as to a
King, and to God, and to One Who is dead for thee.''

\subsubsection*{Why he was bound}

The Fathers dwell on the binding because they know why it was done. St Ambrose
writes, \latin{ille involutus in pannis ut tu mortis laqueis absolutus}: ``he was
wrapped in swaddling bands, that you might be loosed from the snares of
death.'' A few lines further on he says what a Christian can say before the crib:

\begin{quote}
I owe more, O Lord Jesus, to Thy sufferings that I am redeemed, than to Thy
works that I am created. It were no advantage to be born, had it not advantaged
me to be redeemed also.\par
{\smallskip\noindent\small ---St Ambrose, On St Luke, Book 2\par}
\end{quote}

The bands did not hold him. St Gregory Nazianzen sets each lowliness of the
Nativity beside its glory:

\begin{quote}
He was wrapped in swaddling clothes---but He took off the swathing bands of the
grave by His rising again. He was laid in a manger---but He was glorified by
Angels, and proclaimed by a star, and worshipped by the Magi.\par
{\smallskip\noindent\small ---St Gregory Nazianzen, Third Theological Oration\par}
\end{quote}

On Easter morning St Peter went into the tomb:

\begin{quote}
Then cometh Simon Peter, following him, and went into the sepulchre: and saw
the linen cloths lying, and the napkin that had been about his head, not lying
with the linen cloths, but apart, wrapped up into one place.\par
{\smallskip\noindent\small ---John 20:6--7\par}
\end{quote}

The bands of the tomb were left lying, and the one they had bound was risen.

On 19 December the Church sang to him as the Root of Jesse, in an antiphon made
from a verse of Isaiah. The antiphon stops before the end of the verse. Isaiah
goes on, and his last words may be read here, beside the empty tomb:

\begin{quote}
In that day the root of Jesse, who standeth for an ensign of the people, him
the Gentiles shall beseech, and his sepulchre shall be glorious.\par
{\smallskip\noindent\small ---Isaiah 11:10\par}
\end{quote}

On one evening before the feast, read the second chapter of St Luke as far as the
angels' song, and then the burial in St John's nineteenth chapter, from the
coming of Joseph of Arimathea to the end. Then look at the crib. The Child is
wrapped and laid for you.
